% ============================================================
%  BANK SOAL MATEMATIKA SMA — KURIKULUM MERDEKA
%  BAGIAN 3 — FASE F+  |  BLOK 1 (No. 101–128)  |  VERSI SOAL
%  Bilangan Kompleks, Polinomial, Matriks
%  - Tanpa banner; opsi A-E tersusun ke bawah
%  - Sinkron dengan Versi Lengkap (soal & nomor identik)
%  Kelas dokumen: exam  |  Kompilasi: pdfLaTeX (Overleaf)
% ============================================================
\documentclass[11pt,a4paper]{exam}

\usepackage[utf8]{inputenc}
\usepackage[T1]{fontenc}
\usepackage[bahasa]{babel}
\usepackage{amsmath,amssymb}
\usepackage{geometry}
\usepackage{xcolor}
\usepackage{tikz}
\usetikzlibrary{arrows.meta}
\usepackage{pgfplots}
\pgfplotsset{compat=1.18}
\usepackage{enumitem}
\usepackage{array,booktabs}

\geometry{a4paper,margin=2.3cm}

% (Versi Soal: TANPA \printanswers)

\definecolor{biru}{HTML}{1F4E79}
\definecolor{hijau}{HTML}{2E7D32}
\definecolor{oranye}{HTML}{C75B12}
\definecolor{ungu}{HTML}{6A1B9A}

\pagestyle{headandfoot}
\runningheadrule
\firstpageheadrule
\header{\textbf{\textcolor{biru}{BANK SOAL MATEMATIKA}}}{}{\textbf{FASE F+ · Blok 1}}
\footer{}{Halaman \thepage}{}

\renewcommand{\choicelabel}{\Alph{choice}.}

\renewcommand{\choiceshook}{%
  \setlength{\leftmargin}{2em}%
  \setlength{\itemsep}{5pt}%
  \setlength{\topsep}{5pt}%
  \setlength{\parsep}{0pt}}

\newcommand{\Fund}{{\footnotesize\textsf{\textcolor{hijau}{[Fundamental]}}}\ }
\newcommand{\Inter}{{\footnotesize\textsf{\textcolor{biru}{[Intermediate]}}}\ }
\newcommand{\Adv}{{\footnotesize\textsf{\textcolor{oranye}{[Advanced]}}}\ }
\newcommand{\Exp}{{\footnotesize\textsf{\textcolor{ungu}{[Expert]}}}\ }

\newcommand{\topik}[1]{%
  \par\addvspace{14pt}%
  \noindent{\large\bfseries\color{biru}#1}\par\addvspace{6pt}}

\begin{document}

\begin{center}
{\color{biru}\rule{\linewidth}{1.5pt}}\\[4pt]
{\LARGE\bfseries\color{biru} BANK SOAL MATEMATIKA SMA}\\[2pt]
{\large Versi Soal — Fase F+ Blok 1}\\[2pt]
{\large Bilangan Kompleks · Polinomial · Matriks (No. 101–128)}\\[4pt]
{\color{biru}\rule{\linewidth}{1.5pt}}
\end{center}
\vspace{4pt}

\noindent\textit{Petunjuk:} Pilih satu jawaban yang paling tepat. Label tingkat
kesulitan tertera pada awal tiap soal.

\setcounter{question}{100}

\topik{13. Bilangan Kompleks (8 Soal)}

\begin{questions}

\question \Fund Nilai $i^{2}$ adalah \dots
\begin{choices}
\choice $1$ \choice $i$ \CorrectChoice $-1$ \choice $-i$ \choice $0$
\end{choices}

\question \Fund Hasil $(3+2i)+(1-4i)$ adalah \dots
\begin{choices}
\choice $4+2i$ \choice $2-2i$ \choice $4-6i$ \CorrectChoice $4-2i$ \choice $2+6i$
\end{choices}

\question \Inter Hasil $(2+i)(3-i)$ adalah \dots
\begin{choices}
\choice $6+i$ \choice $5+i$ \CorrectChoice $7+i$ \choice $7-i$ \choice $6-i$
\end{choices}

\question \Inter Modulus dari $z = 3+4i$ adalah \dots
\begin{choices}
\choice $7$ \CorrectChoice $5$ \choice $25$ \choice $\sqrt{7}$ \choice $12$
\end{choices}

\question \Inter Konjugat dari $z = 5-2i$ adalah \dots
\begin{choices}
\choice $-5+2i$ \choice $2-5i$ \choice $-5-2i$ \CorrectChoice $5+2i$ \choice $5-2i$
\end{choices}

\question \Adv Pada diagram Argand berikut, bilangan kompleks $z$ digambarkan
sebagai vektor.

\begin{center}
\begin{tikzpicture}[scale=0.9]
\coordinate (O) at (0,0);
\coordinate (Z) at (1.5,1.5);
\draw[->] (-0.3,0)--(2.3,0) node[right]{Re};
\draw[->] (0,-0.3)--(0,2.3) node[above]{Im};
\draw[very thick,biru,-{Latex}] (O)--(Z) node[above right]{$z=1+i$};
\draw[dashed,gray] (1.5,0)--(Z) (0,1.5)--(Z);
\node[below] at (1.5,0){\small $1$}; \node[left] at (0,1.5){\small $1$};
\end{tikzpicture}
\end{center}

Bentuk polar dari $z$ adalah \dots
\begin{choices}
\choice $2(\cos 45^\circ + i\sin 45^\circ)$
\CorrectChoice $\sqrt{2}(\cos 45^\circ + i\sin 45^\circ)$
\choice $\sqrt{2}(\cos 90^\circ + i\sin 90^\circ)$
\choice $1(\cos 45^\circ + i\sin 45^\circ)$
\choice $\sqrt{2}(\cos 60^\circ + i\sin 60^\circ)$
\end{choices}

\question \Adv Nilai $i^{2023}$ adalah \dots
\begin{choices}
\choice $i$ \choice $1$ \choice $-1$ \CorrectChoice $-i$ \choice $0$
\end{choices}

\question \Exp Modulus dan argumen dari $z = -1+\sqrt{3}\,i$ adalah \dots
\begin{choices}
\choice modulus $2$, argumen $60^\circ$
\choice modulus $4$, argumen $120^\circ$
\CorrectChoice modulus $2$, argumen $120^\circ$
\choice modulus $2$, argumen $240^\circ$
\choice modulus $\sqrt{2}$, argumen $120^\circ$
\end{choices}

\topik{14. Polinomial (10 Soal)}

\question \Fund Derajat polinomial $3x^{4} - 2x^{2} + 5$ adalah \dots
\begin{choices}
\choice $5$ \CorrectChoice $4$ \choice $3$ \choice $2$ \choice $1$
\end{choices}

\question \Fund Sisa pembagian $P(x) = x^{2}+2x+1$ oleh $(x-1)$ adalah \dots
\begin{choices}
\choice $1$ \choice $0$ \CorrectChoice $4$ \choice $3$ \choice $5$
\end{choices}

\question \Inter Sisa pembagian $P(x) = 2x^{3}+x^{2}-7$ oleh $(x-2)$ adalah \dots
\begin{choices}
\choice $9$ \CorrectChoice $13$ \choice $11$ \choice $7$ \choice $15$
\end{choices}

\question \Inter Faktor $(x-2)$ merupakan faktor dari $P(x)$ jika \dots
\begin{choices}
\choice $P(0)=2$ \choice $P(-2)=0$ \CorrectChoice $P(2)=0$ \choice $P(2)=2$ \choice $P'(2)=0$
\end{choices}

\question \Inter Tabel berikut memuat beberapa nilai $P(x) = x^{3}-2x^{2}-5x+6$.

\begin{center}\small
\begin{tabular}{|c|c|c|c|c|c|}
\hline
$x$ & $-2$ & $-1$ & $1$ & $2$ & $3$\\\hline
$P(x)$ & $0$ & $8$ & $0$ & $-4$ & $0$\\\hline
\end{tabular}
\end{center}

Berdasarkan tabel, salah satu faktor $P(x)$ adalah \dots
\begin{choices}
\choice $(x+1)$ \CorrectChoice $(x-1)$ \choice $(x-2)$ \choice $(x+3)$ \choice $(x-6)$
\end{choices}

\question \Inter Akar-akar dari $x^{3}-2x^{2}-5x+6 = 0$ adalah \dots
\begin{choices}
\choice $1,\ -3,\ 2$ \choice $-1,\ 3,\ 2$ \CorrectChoice $1,\ 3,\ -2$
\choice $1,\ 2,\ 3$ \choice $-1,\ -3,\ -2$
\end{choices}

\question \Adv Hasil bagi $x^{3}-2x^{2}-5x+6$ oleh $(x-1)$ menggunakan skema Horner
adalah \dots
\begin{choices}
\CorrectChoice $x^{2}-x-6$ \choice $x^{2}+x-6$ \choice $x^{2}-3x+6$
\choice $x^{2}-x+6$ \choice $x^{2}-2x-6$
\end{choices}

\question \Adv Jumlah akar-akar persamaan $2x^{3}-6x^{2}+4x-1 = 0$ adalah \dots
\begin{choices}
\choice $-3$ \choice $2$ \choice $6$ \CorrectChoice $3$ \choice $1$
\end{choices}

\question \Adv Suatu polinomial $P(x)$ dibagi $(x-1)$ bersisa 3 dan dibagi $(x-2)$
bersisa 5. Sisa pembagian $P(x)$ oleh $(x-1)(x-2)$ adalah \dots
\begin{choices}
\choice $2x-1$ \choice $x+2$ \choice $2x+3$ \CorrectChoice $2x+1$ \choice $x+1$
\end{choices}

\question \Exp Jika $x = 2$ merupakan akar dari $x^{3}-3x^{2}+kx-4 = 0$, maka nilai
$k$ adalah \dots
\begin{choices}
\choice $2$ \choice $6$ \CorrectChoice $4$ \choice $0$ \choice $8$
\end{choices}

\topik{15. Matriks (10 Soal)}

\question \Fund Determinan matriks $\begin{pmatrix} 2 & 1\\ 3 & 4 \end{pmatrix}$ adalah \dots
\begin{choices}
\CorrectChoice $5$ \choice $11$ \choice $8$ \choice $-5$ \choice $7$
\end{choices}

\question \Fund Hasil $\begin{pmatrix} 1 & 2\\ 3 & 4 \end{pmatrix} +
\begin{pmatrix} 2 & 0\\ 1 & 5 \end{pmatrix}$ adalah \dots
\begin{choices}
\choice $\begin{pmatrix} 3 & 2\\ 4 & 8 \end{pmatrix}$
\CorrectChoice $\begin{pmatrix} 3 & 2\\ 4 & 9 \end{pmatrix}$
\choice $\begin{pmatrix} 2 & 0\\ 3 & 20 \end{pmatrix}$
\choice $\begin{pmatrix} 3 & 2\\ 3 & 9 \end{pmatrix}$
\choice $\begin{pmatrix} 1 & 2\\ 4 & 9 \end{pmatrix}$
\end{choices}

\question \Inter Invers matriks $\begin{pmatrix} 2 & 1\\ 3 & 4 \end{pmatrix}$ adalah \dots
\begin{choices}
\choice $\dfrac{1}{5}\begin{pmatrix} 2 & -1\\ -3 & 4 \end{pmatrix}$
\CorrectChoice $\dfrac{1}{5}\begin{pmatrix} 4 & -1\\ -3 & 2 \end{pmatrix}$
\choice $\dfrac{1}{5}\begin{pmatrix} 4 & 1\\ 3 & 2 \end{pmatrix}$
\choice $\begin{pmatrix} 4 & -1\\ -3 & 2 \end{pmatrix}$
\choice $\dfrac{1}{11}\begin{pmatrix} 4 & -1\\ -3 & 2 \end{pmatrix}$
\end{choices}

\question \Inter Hasil $2\begin{pmatrix} 1 & -2\\ 0 & 3 \end{pmatrix}$ adalah \dots
\begin{choices}
\CorrectChoice $\begin{pmatrix} 2 & -4\\ 0 & 6 \end{pmatrix}$
\choice $\begin{pmatrix} 2 & -2\\ 0 & 3 \end{pmatrix}$
\choice $\begin{pmatrix} 3 & 0\\ 2 & 5 \end{pmatrix}$
\choice $\begin{pmatrix} 1 & -4\\ 0 & 6 \end{pmatrix}$
\choice $\begin{pmatrix} 2 & -4\\ 0 & 3 \end{pmatrix}$
\end{choices}

\question \Inter Hasil $\begin{pmatrix} 1 & 2\\ 3 & 4 \end{pmatrix}
\begin{pmatrix} 1 & 0\\ 0 & 1 \end{pmatrix}$ adalah \dots
\begin{choices}
\choice $\begin{pmatrix} 1 & 0\\ 0 & 1 \end{pmatrix}$
\CorrectChoice $\begin{pmatrix} 1 & 2\\ 3 & 4 \end{pmatrix}$
\choice $\begin{pmatrix} 1 & 2\\ 0 & 4 \end{pmatrix}$
\choice $\begin{pmatrix} 4 & 3\\ 2 & 1 \end{pmatrix}$
\choice $\begin{pmatrix} 2 & 4\\ 6 & 8 \end{pmatrix}$
\end{choices}

\question \Inter Determinan matriks $\begin{pmatrix} 3 & 2\\ 1 & 4 \end{pmatrix}$ adalah \dots
\begin{choices}
\CorrectChoice $10$ \choice $14$ \choice $11$ \choice $-10$ \choice $12$
\end{choices}

\question \Adv Di sebuah toko, harga 2 buku dan 1 pensil Rp5.000, sedangkan 3 buku
dan 4 pensil Rp10.000. Dengan metode matriks, harga satu buku adalah \dots
\begin{choices}
\CorrectChoice Rp2.000 \choice Rp1.000 \choice Rp3.000 \choice Rp4.000 \choice Rp2.500
\end{choices}

\question \Adv Matriks $\begin{pmatrix} 2 & k\\ 3 & 6 \end{pmatrix}$ tidak memiliki
invers (singular) jika $k = \dots$
\begin{choices}
\choice $2$ \choice $3$ \choice $6$ \CorrectChoice $4$ \choice $0$
\end{choices}

\question \Adv Jika $A = \begin{pmatrix} 1 & 2\\ 3 & 4 \end{pmatrix}$, maka $\det(2A)$ adalah \dots
\begin{choices}
\choice $-2$ \choice $-4$ \choice $8$ \CorrectChoice $-8$ \choice $-16$
\end{choices}

\question \Exp Jika $A = \begin{pmatrix} 1 & 1\\ 0 & 1 \end{pmatrix}$, maka $A^{3}$ adalah \dots
\begin{choices}
\choice $\begin{pmatrix} 1 & 1\\ 0 & 1 \end{pmatrix}$
\choice $\begin{pmatrix} 3 & 3\\ 0 & 3 \end{pmatrix}$
\CorrectChoice $\begin{pmatrix} 1 & 3\\ 0 & 1 \end{pmatrix}$
\choice $\begin{pmatrix} 1 & 3\\ 0 & 3 \end{pmatrix}$
\choice $\begin{pmatrix} 3 & 1\\ 0 & 1 \end{pmatrix}$
\end{choices}

\end{questions}

\vspace{6pt}
\begin{center}
{\color{biru}\rule{\linewidth}{1pt}}\\[3pt]
{\small\textit{Akhir F+ Blok 1 (101–128) — Versi Soal. Kunci \& pembahasan ada pada Versi Lengkap.}}
\end{center}

\end{document}